Pour chacune des affirmations suivantes, indiquer si elle est VRAIE ou FAUSSE en justifiant la réponse.

\textit{Toutes les affirmations sont indépendantes.}

\begin{enumerate}
\item On considère la suite $(u_n)$ définie sur $\mathbb{N}$ par $u_0=5$ et $u_{n+1}=2u_n+1$.

\textbf{Affirmation 1 :} $u_3=47$.

\bigskip

\item Mathétix anime une chaîne vidéo sur le thème des mathématiques.

Il observe qu'entre deux mois consécutifs, il perd un dixième de ses abonnés mais il en gagne 200.

On modélise cette évolution par la suite $(V_n)$, avec $V_n$ le nombre d'abonnés le $n$-ième mois.

\textbf{Affirmation 2 :} On a : $V_{n+1}=1{,}1V_n+200$ pour tout entier naturel $n$.

\bigskip

\item On considère la suite $(w_n)$ définie sur $\mathbb{N}$ par $w_n=\dfrac{n+4}{3}$.

\textbf{Affirmation 3 :} La suite $(w_n)$ est une suite arithmétique de premier terme 4 et de raison $\dfrac{1}{3}$.

\bigskip

\item On interroge 200 élèves sur leur moyen de transport pour venir au lycée.
\begin{center}
\begin{tabular}{|c|c|c|c|}
\hline
 & \textbf{Filles} & \textbf{Garçons} & \textbf{Total}\\
\hline
\textbf{Bus} & 90 &  & 150\\
\hline
\textbf{Autre} & 15 & 35 & 50\\
\hline
\textbf{Total} & 105 & 95 & 200\\
\hline
\end{tabular}
\end{center}
On utilisera ce tableau pour les affirmations 4 \underline{et} 5.

On choisit un élève au hasard.

\bigskip

\textbf{Affirmation 4 :} La probabilité qu'il se rende au lycée en bus est de $\dfrac{3}{4}$.

\bigskip

\textbf{Affirmation 5 :} Sachant que l'élève a pris le bus, la probabilité que ce soit un garçon est de $0{,}6$.
\end{enumerate}
